% ============================================================
%  MO DAU
% ============================================================
\unnumberedchapter{MỞ ĐẦU}

\unnumberedsection{1. Lý do chọn đề tài}

Bữa cơm gia đình là một hoạt động diễn ra hằng ngày và liên tục. Đi kèm với nó
là một chuỗi quyết định nhỏ nhưng lặp lại: nấu món gì, cần mua thêm nguyên liệu
nào, làm sao để cả tuần không bị trùng món và vẫn cân đối về dinh dưỡng. Người
đảm nhiệm việc bếp núc trong gia đình phải xử lý chuỗi quyết định này bằng trí
nhớ và kinh nghiệm cá nhân.

Thực tế đó dẫn tới ba khó khăn cụ thể.

\textbf{Thứ nhất, khó tận dụng nguyên liệu sẵn có.} Người dùng thường ở trong
tình huống ngược với các công cụ hiện có: họ không tìm kiếm theo tên món, mà
xuất phát từ những nguyên liệu đang có trong tủ lạnh và cần biết có thể nấu được
món gì. Việc tra cứu theo chiều này đòi hỏi phải đối chiếu thủ công từng công
thức, một việc gần như không khả thi khi số lượng công thức lớn.

\textbf{Thứ hai, khó lập thực đơn cho cả tuần.} Lập thực đơn tuần là bài toán có
ràng buộc: phải tránh lặp món, phải phù hợp với từng bữa trong ngày, và nên cân
đối năng lượng. Khi làm thủ công, người dùng có xu hướng quay lại một nhóm nhỏ
các món quen thuộc, làm bữa ăn trở nên đơn điệu.

\textbf{Thứ ba, khó lập danh sách đi chợ chính xác.} Sau khi đã có thực đơn,
việc tổng hợp nguyên liệu cần mua đòi hỏi phải cộng dồn khối lượng của cùng một
nguyên liệu xuất hiện ở nhiều món khác nhau. Sai sót ở bước này dẫn tới mua
thiếu, phải đi chợ nhiều lần, hoặc mua thừa gây lãng phí.

Ba khó khăn trên có bản chất là các bài toán xử lý dữ liệu có cấu trúc, hoàn
toàn có thể tự động hóa bằng phần mềm. Đây chính là lý do đề tài được lựa chọn:
xây dựng một hệ thống giải quyết trọn vẹn cả ba khâu trong một luồng công việc
liền mạch, thay vì chỉ giải quyết riêng lẻ từng khâu.

\unnumberedsection{2. Mục tiêu của đề tài}

Đồ án đặt ra các mục tiêu cụ thể sau:

Một là xây dựng ứng dụng web quản lý được kho dữ liệu món ăn Việt Nam gồm công
thức, nguyên liệu, danh mục, vùng miền, độ khó và thông tin năng lượng. Hai là
thiết kế và cài đặt thuật toán \textbf{gợi ý món ăn theo tập nguyên liệu sẵn có}
biết xếp hạng kết quả theo mức độ phù hợp và chỉ rõ nguyên liệu còn thiếu của
từng món. Ba là thuật toán \textbf{sinh thực đơn tuần tự động} với các ràng buộc
về tính đa dạng, sự phù hợp theo bữa và cân đối năng lượng, đồng thời cho phép
chỉnh sửa thủ công từng suất ăn. Bốn là chức năng \textbf{sinh danh sách đi chợ}
tự động từ thực đơn đã lập, có gộp nhóm và cộng dồn khối lượng. Năm là cơ chế
xác thực và phân quyền theo vai trò, tách bạch chức năng của khách, người dùng
đã đăng nhập và quản trị viên. Sáu là kiểm chứng tính đúng đắn của các thuật
toán lõi bằng kiểm thử đơn vị tự động.

\unnumberedsection{3. Đối tượng và phạm vi nghiên cứu}

\textbf{Đối tượng nghiên cứu.} Đối tượng của đề tài gồm hai nhóm. Về mặt người
dùng, đó là người trực tiếp nấu ăn trong gia đình Việt Nam, có nhu cầu lên thực
đơn và đi chợ định kỳ. Về mặt kỹ thuật, đó là các phương pháp biểu diễn quan hệ
nhiều - nhiều giữa món ăn và nguyên liệu, các độ đo tương đồng tập hợp phục vụ
xếp hạng, và các thuật toán tham lam có ràng buộc phục vụ bài toán xếp lịch.

\textbf{Phạm vi nghiên cứu.} Đồ án giới hạn trong các nội dung sau:

Hình thức triển khai là \textbf{ứng dụng web} chạy trên trình duyệt, không phát
triển ứng dụng di động gốc. Thuật toán gợi ý dựa trên \textbf{đối sánh tập hợp
nguyên liệu}, không dùng học máy hay lọc cộng tác, vì hệ thống mới chưa có dữ
liệu hành vi người dùng để huấn luyện mô hình. Thuật toán lập thực đơn theo
hướng \textbf{tham lam có ràng buộc}, không giải bài toán tối ưu toàn cục. Thông
tin dinh dưỡng giới hạn ở mức \textbf{năng lượng (kcal) theo khẩu phần}. Dữ liệu
mẫu gồm 26 món ăn Việt Nam phổ biến, đủ để minh họa và kiểm chứng thuật toán,
không nhằm xây dựng kho công thức quy mô lớn.

\unnumberedsection{4. Phương pháp nghiên cứu}

Đồ án kết hợp bốn phương pháp: \textbf{khảo sát} các ứng dụng nấu ăn và lập thực
đơn hiện có trong và ngoài nước để xác định những chức năng đã được đáp ứng và
những khoảng trống còn lại; \textbf{phân tích và thiết kế hệ thống} bằng sơ đồ
use case, thiết kế cơ sở dữ liệu quan hệ, thiết kế lớp và sơ đồ tuần tự;
\textbf{thực nghiệm} bằng cách cài đặt hệ thống, chạy thử trên dữ liệu mẫu và
quan sát kết quả đầu ra của các thuật toán; và \textbf{kiểm thử} bằng kiểm thử
đơn vị tự động cho các thuật toán lõi, dùng kiểm thử để phát hiện khiếm khuyết
và xác nhận hiệu quả sau khi sửa.

\unnumberedsection{5. Bố cục báo cáo}

Ngoài phần Mở đầu và phần Kết luận, báo cáo được tổ chức thành năm chương.
\textbf{Chương 1} trình bày bối cảnh thực tiễn, khảo sát các ứng dụng tương tự
và xác định vấn đề mà đồ án tập trung giải quyết. \textbf{Chương 2} trình bày cơ
sở lý thuyết về kiến trúc MVC, ASP.NET Core, Entity Framework Core, ASP.NET Core
Identity, SQL Server và nền tảng thuật toán được sử dụng. \textbf{Chương 3} mô
tả bài toán, đặc tả yêu cầu, thiết kế cơ sở dữ liệu, lược đồ use case, kiến trúc
hệ thống, thiết kế lớp, thiết kế luồng xử lý và chi tiết cài đặt các thuật toán.
\textbf{Chương 4} giới thiệu giao diện các chức năng, kịch bản demo và kết quả
kiểm thử. \textbf{Chương 5} nêu kết luận và hướng phát triển.
